\section{Discussion}
\label{sec:discussion}

\paragraph{The size of the cost depends on the step and on the sample.} The results say that the holdout of the protocol is an unbiased estimate of the fitted model's AUC, and that breaking the rules does not cost the same everywhere. A step that reads only inputs costs nothing measurable in the estimate (twelve of twelve cells); a step that reads labels costs between \nv{syn.w.scr.bias} (the screen) and \nv{syn.w.fit.bias} (every step) on 2,500 planted rows, and about twice as much on 800 rows; a calibration map fitted on the holdout costs about as much as the weight of evidence with holdout labels, and unlike the other leaks it makes the calibration error look better than it is. The ordering follows the flexibility of the step. A weight of evidence has ten free values per input, a screen picks ten of 40 inputs by their association with the label, a scaling has two parameters per input. The ordering is consistent with the optimism of an in-sample statistic, which grows with the number of quantities fitted from the rows that the estimate later reads, and the dependence on the sample size in the conjecture of Section~\ref{sec:theory} is the same mechanism seen through the cross-validation.

\paragraph{Why the interval misleads more than the point estimate.} The coverage of the engine's interval falls to \nv{syn.w.fit.cov}\% and, for the product's own scorecard, to \nv{prod.p.fit.cov}\% when every step reads the holdout. The interval is the DeLong interval, which describes the sampling variability of the AUC of a fixed model on a holdout. A leak does not change that variability; it shifts the centre by an amount of the order of the interval's half-width, so that the interval is a statement about the sampling error of a number that is already biased. A validator who reads the interval as the range in which the future AUC lies gets a range that excludes the future most of the time.

\paragraph{What the gates tell a validator and what they do not.} The gates fail whenever a rule is broken and they fail with the same certainty whatever the cost of the breach. The leak through imputation on all rows costs nothing in the estimate and fails two gates in every draw. For a validator the useful reading is that a failure localises the breach (a transform, the calibration map or the choice among configurations) and that the bias table says what that kind of breach has cost on data of the same shape. A protocol that reported only the cost would let a harmless leak pass, and a pipeline that is harmless on one sample can be harmful on another, as the doubling of the effects on the small sample shows. The gate that fails on a harmless breach is the price of a test that does not depend on the data.

\paragraph{Discrete choices need more than one draw.} The only leak that the third gate does not catch in every draw at its default setting is the choice of a configuration on the holdout, and Proposition~\ref{prop:discrete} says why: with two configurations a perturbed choice agrees with the base choice half of the time. The measurement agrees with the proposition (\nv{e2.draws.l.1}\% with one draw for the logit and \nv{e2.draws.l.5}\% with five), and the practical rule is that a pipeline that chooses among $K$ configurations needs $d$ draws with $K^{-d}$ below the miss probability that the validator accepts: five draws bring the detection of a choice among two configurations to \nv{e2.draws.pred.l.5}\% by the proposition, and two draws bring that of a choice among eight to \nv{e2.draws.pred.w.2}\%. The second gate catches the same leak in \nv{e2.w.sel.G2}\% and \nv{e2.l.sel.G2}\% of the draws through the calibration map, so that the protocol with all three gates detects it in every draw, but the dependence of that detection on the map is an accident of the pipeline and not a guarantee.

\paragraph{What the engine's tests see of a leak that exists only in development.} The planted leak shows a division of labour that a validation function should keep in mind. A test that compares the current AUC with the development AUC sees the leak at a shift of half a standard deviation (\nv{e3.s0.5.auc_test}\%); the stability index of the leaked feature sees it \nv{e3.s1.0.csi_test}\% of the time at a shift of one, and the single-feature screen not at all. A monitoring process that tracks input drift would therefore report a stable input and a falling AUC, so that a stable input does not certify the absence of a feature that carries the outcome in development and not at decision time. Proposition~\ref{prop:psi} says that this is not an accident of the generator: the index of a feature that is informative for a small share of the rows is second order in that share. The consequence for the practice that the protocol serves is that the leak that the gates cannot reach, a feature built from information dated after the observation date, has to be sought where it shows, in the discrimination out of time, and that a model whose development AUC is much higher than its AUC on the first cohort that follows is suspect before it is drifted.

\paragraph{A declared cohort measures the level, and a random split does not.} On the loan data a random sample of the training window gives the right level of the 5-year rate in every cohort, within \nv{e4.holdse} points, and the declared out-of-time cohort gives gaps of up to \nv{e4.2006-07.sc.aoot_gap} points. The AUC does not show the same pattern, and a random split is not systematically optimistic in discrimination on these data, which is consistent with a rank order that is more stable across vintages than the level of the default rate. For the provisions and the capital that use the PD as a level, the split that matters is the declared cohort, and a case that has none, as the card holders and the Polish bankruptcy files do, states that it makes no out-of-time claim instead of leaving the reader to assume it.

\paragraph{Negative and limiting results.} One prediction did not hold: the screen with holdout labels does not raise the AUC on the card holders, so the claim that every step that reads labels inflates the estimate is not established for that step on that data set. The leaks through inputs only do not inflate the estimate in any of the twelve cells, so the cost of a leak cannot be inferred from the fact of the breach. The choice among two configurations on the holdout costs nothing measurable. And a closed form that was expected to be accurate for the continuous index (Proposition~\ref{prop:psi}) is accurate only to a shift of one half; the binned form is the one to use. These limits narrow the claim of the paper to the one the results sustain: with the rules enforced and tested, the holdout number is unbiased and its interval covers; without them, the cost depends on the step and the sample, and the gates detect the breach whatever its cost.

\paragraph{Where the protocol sits in the practice of a validation function.} The protocol is a part of what a validator does before reading a holdout AUC. It is cheap ($4+d$ fits of the pipeline), it needs no change of the model and it yields a result that can be written to the artifact of a case, with the split and its limits. It is not a substitute for the review of the features for information from the future, for the out-of-time cohort where the data allow one, or for the discrimination test against the development value that the supervisory instructions ask for. The results of Section~\ref{sec:results} show that each of these sees something that the others do not.
